\documentclass[11pt]{article}
\def\activityid{F10-F-v2}\usepackage[letterpaper,margin=.65in,top=.60in,bottom=.65in]{geometry}
\usepackage[T1]{fontenc}
\usepackage[default]{sourcesanspro}
\usepackage{microtype,tikz,array,tabularx,fancyhdr,amsmath,amssymb}
\usepackage[hidelinks]{hyperref}
\usetikzlibrary{calc,arrows.meta}
\setlength{\parindent}{0pt}\setlength{\parskip}{7pt}
\setlength{\headheight}{0pt}\setlength{\footskip}{26pt}\setlength{\emergencystretch}{2em}
\pagestyle{fancy}\fancyhf{}\renewcommand{\headrulewidth}{0pt}\renewcommand{\footrulewidth}{.4pt}
\fancyfoot[L]{\scriptsize Bellingham Math Circle / Week 10 / \activityid}
\fancyfoot[R]{\scriptsize \thepage}
\definecolor{ink}{gray}{.12}\definecolor{soft}{gray}{.95}\color{ink}
\newcommand{\PageTitle}[2]{{\fontsize{10}{12}\selectfont\bfseries WEEK 10\quad ROUTE LAB\hfill #1\par}\vspace{3pt}{\fontsize{26}{29}\selectfont\bfseries #2\par}\vspace{2pt}}
\newcommand{\NameLine}{{\small Name: \rule{2.65in}{.4pt}\hfill Date: \rule{1.35in}{.4pt}\par}\vspace{4pt}}
\newcommand{\Task}[2]{\par\vspace{3pt}{\large\bfseries #1}\par #2\par}
\newcommand{\Subhead}[1]{\par\vspace{3pt}{\large\bfseries #1}\par}
\newcommand{\RuleBox}[1]{\par\begingroup\setlength{\fboxsep}{8pt}\colorbox{soft}{\parbox{\dimexpr\linewidth-16pt\relax}{#1}}\endgroup\par}
\newcommand{\WriteBox}[2]{\par\textbf{#1}\par\vspace{2pt}\begin{tikzpicture}\draw[black!40,line width=.5pt] (0,0) rectangle (\dimexpr\linewidth-.5pt\relax,#2);\end{tikzpicture}\par}
\newcommand{\StudentFooter}{\vfill{\small Try it with objects. Explain by pointing, talking, or drawing.}\par}
\newcommand{\RookBoard}[3]{\begin{tikzpicture}[x=#3,y=#3]\draw[step=1,black!45] (0,0) grid (#1,#2);\draw[thick] (0,0) rectangle (#1,#2);\node[font=\Large] at ({#1-.5},.5){$\star$};\draw[-{Latex},thick] (0,-.35)--(#1,-.35);\draw[-{Latex},thick] ({#1+.35},#2)--({#1+.35},0);\end{tikzpicture}}
\newcommand{\Table}[3]{\begin{tikzpicture}[x=#3,y=#3]\draw[step=1,black!25] (0,0) grid (#1,#2);\draw[thick] (0,0) rectangle (#1,#2);\fill (0,0)circle(3pt);\draw[-{Latex},very thick] (0,0)--(.7,.7);\node[below] at ({#1/2},-.1){width #1};\node[rotate=90,above] at (-.2,{#2/2}){height #2};\end{tikzpicture}}
\newcommand{\GraphNodes}{\coordinate(A)at(0,0);\coordinate(B)at(3,0);\coordinate(C)at(1.5,2.6);\coordinate(D)at(1.5,.95);}
\newcommand{\Kfour}[1]{\begin{tikzpicture}[scale=#1]\GraphNodes\draw[very thick](A)--(B)--(C)--(A)--(D)--(B) (D)--(C);\foreach \n in {A,B,C,D}{\node[circle,draw,fill=white,minimum size=22pt,font=\bfseries]at(\n){\n};}\end{tikzpicture}}
\newcommand{\House}[1]{\begin{tikzpicture}[scale=#1]\coordinate(A)at(0,0);\coordinate(B)at(3,0);\coordinate(C)at(3,2);\coordinate(D)at(0,2);\coordinate(E)at(1.5,3.3);\draw[very thick](A)--(B)--(C)--(D)--(A) (D)--(E)--(C);\foreach \n in {A,B,C,D,E}{\node[circle,draw,fill=white,minimum size=22pt,font=\bfseries]at(\n){\n};}\end{tikzpicture}}

\begin{document}

\PageTitle{FACILITATOR / 1 OF 4}{Puzzle makers become route planners}
\RuleBox{Core question: When is an every-road-once walk possible, and when repetition is necessary, how little can we repeat? Begin by moving a counter and marking used roads. Every drawing has a connected network and undirected roads. Each road joins two different dots; parallel roads are allowed.}
\Subhead{Preparation and access}
Print 3 K--1, 3 middle, 1 upper, and one extra reserve page. Bring seven moving counters, about 35 small counters to mark used edges, pencils, scrap paper, and two short paper strips for the upper loop splice. Strings between paper islands are an optional larger model. Setup takes about 10 minutes. No ruler accuracy or actual bridge construction is required.

K--1: no reading; distinguish a line from its endpoints, keep track of used bridges, reason about dead ends. Middle: count degrees to 4 and recognize odd/even; understand that failure of a few attempts is not proof. Upper: follow a route list, count edges, give a general argument with paired arrivals/departures. Extra: add lengths to 23, compare all three pairings, follow a shortest-path and lower-bound argument; no algebra beyond these small sums.
\Subhead{A shared hour with a meaningful workshop}
0--10 build little maps and walk freely. 10--15 briefly demonstrate every-bridge-once, allowing island revisits. 15--35 main investigation: parent starts K--1; organizer launches the house and two-triangle splice, then alternates visits about every five minutes. 35--40 movement/reset. 40--55 continue repairs or shortest delivery, or choose a proof continuation. 55--60 share a reason and tidy.

Triplets rotate walker, edge-marker, and skeptic. Parent stays mostly with K--1; organizer visits middle and upper about every five minutes. The fifth grader should test a conjecture against an adult and explain a lower bound, not merely receive harder numbers.
\Subhead{Launch and hint sequence}
``Can you cross every bridge once without teleporting?'' Initially leave start and finish free. Later ask ``What must happen after you arrive at a place in the middle?'' Pair incident edges with fingers. Reveal odd/even vocabulary after the physical pairing. For route optimization ask why repeating one road changes parity at only two places. Offer the three possible pairings only after a first route attempt.

The puzzle-making ending has a concrete specification, a construction, and a certificate. Archive the route as well as the picture; testing by one partner alone does not prove optimality or impossibility.

\Subhead{A satisfying stop; an optional proof}
\textbf{K--1:} explain the three-leaf obstruction and show a repair; continue by forcing its endpoints. \textbf{Middle:} explain house endpoints and repair K4; continue with the general odd-endpoint argument. \textbf{Upper:} splice the two triangle loops, then attain and certify delivery length 8; continue with the all-even construction and temporary-edge proof. Record only the claims actually explained.

\newpage

\PageTitle{FACILITATOR / 2 OF 4}{One-walk maps: checks and proof}
\Subhead{K--1 answers}
A triangle has a closed route, e.g. A--B--C--A. Add a branch A--D: a trail is D--A--B--C--A, with endpoints D and A. At a degree-1 leaf, a walk cannot arrive and then leave by a different unused bridge. Thus the three-leaf star cannot have a trail covering every edge once: it would need three endpoints, while a walk has at most two.

Adding a bridge B--C between two leaves repairs the star: D--A--B--C--A is a complete trail. It begins/ends at the remaining leaf D and center A. Accept another pair of leaves with relabeled route.
\Subhead{Middle and upper map answers}
House degrees: A=2, B=2, C=3, D=3, E=2. One trail is D--A--B--C--D--E--C. Every complete trail must start at one of C,D and end at the other. No closed Euler trail exists.

For the triangle with center D, all four vertices A,B,C,D have degree 3. It has no every-edge-once trail. Add a second AB bridge; now A,B are even and C,D odd. A valid trail is C--A--B--D--A--B--C--D, using the two distinct AB bridges separately. Middle children may draw the added bridge as a curve; a crossing without a dot is not a new junction.
\Subhead{Euler criterion with a construction}
In any completed trail, arrivals and departures at an internal vertex pair. A closed Euler trail therefore has only even degrees; an open Euler trail has exactly two odd vertices, its endpoints. The necessity argument includes repeated visits to a vertex: each visit contributes another pair.

First physically splice the two triangles ABC and ADE sharing A. Mark AB, BC, CA used after A--B--C--A; AD, DE, EA remain. Insert the strip A--D--E--A at a visit to A to obtain, for example, A--D--E--A--B--C--A. Trace it and mark all six roads. This one construction motivates the next general argument.

For sufficiency, assume a finite connected graph with at least one edge and all degrees even. Follow unused edges until stuck. At a vertex other than the start, arriving uses one edge of an even available count and leaves at least one to depart; so only the start can be where the walk stops. If unused edges remain, connectedness ensures that some vertex of the current closed walk touches unused edges (take a path toward an unused edge). Start a new closed walk there in the remaining even-degree graph and splice it into the first. Repeat. Finitely many edges guarantee completion.

For exactly two odd vertices, add one temporary edge between them, construct the closed Euler trail, then remove that temporary edge to obtain an open Euler trail. Parallel edges are allowed. This proves existence; it does not promise that every arbitrary greedy walk uses all edges without splicing.

\newpage

\PageTitle{FACILITATOR / 3 OF 4}{Delivery routes and optimality}
\Subhead{Upper: six unit-length roads}
The triangle-and-center network is $K_4$, with six unit roads. Every closed covering walk must repeat at least two road crossings: all four original vertices are odd, and each added traversal changes parity at its two endpoints only. A closed walk has even traversal-degree everywhere. Thus six original plus at least two repeats gives a lower bound of 8.

A matching route is A--B--C--D--A--C--D--B--A. It repeats AB and CD; all other roads occur once. The three possible odd-vertex pairings are AB+CD, AC+BD, AD+BC. Each adds two roads and creates an even connected multigraph, so each admits a closed Euler trail of length 8.

If finishing anywhere is allowed, at most two odd vertices may remain. At least one repeated crossing is necessary, giving a lower bound of 7. A--B--C--A--D--C--D--B attains 7, repeating CD.

For the design brief, a star with five leaves has six odd vertices: five degree-1 leaves and the degree-5 center. But its leaves are not adjacent, so it cannot meet the three-extra-step target. A workable design is two triangles joined by three corresponding cross edges (a triangular prism graph). All six degrees are 3. Duplicate those three cross edges; now all degrees are 4, so a closed tour exists with exactly three extra crossings, meeting the parity lower bound. Label the triangles ABC and DEF, with cross edges AD, BE, CF. One 12-step closed route is A--B--C--A--D--E--B--E--F--C--F--D--A. Children may find other designs.
\Subhead{Extra: six roads of different lengths}
Original total is $2+3+7+1+6+1=20$. Distances between odd vertices are $d(A,B)=2$, $d(A,C)=2$ via D, $d(A,D)=1$, $d(B,C)=3$, $d(B,D)=3$ via A, $d(C,D)=1$.

Pair costs: AB+CD=3; AC+BD=5; AD+BC=4. Choosing the cheapest available AD first is not guaranteed optimal. Duplicate AB and CD to obtain minimum total 23. The same A--B--C--D--A--C--D--B--A route has length 23. Long roads AC and BD must still be serviced once; shortest connections are used only for the \emph{extra} travel.

For a general lower bound, subtract one required copy of every road from a closed covering walk. The remaining multigraph has odd degree exactly at the original odd vertices. Its edges can be decomposed into paths pairing those odd vertices and closed loops: follow unused edges from an odd vertex to another odd one, remove that path, and continue, then remove cycles. Each path costs at least the shortest distance of its endpoint pair, and loops cannot improve a nonnegative-cost total. Thus every solution pays at least the best pairing cost. Duplicating shortest paths for that pairing attains the bound by Euler's construction.

\newpage

\PageTitle{FACILITATOR / 4 OF 4}{The mathematics behind the workshop}
\Subhead{Undergraduate graph theory}
The optional proof conversation completes the Euler trail criterion, including its constructive direction; the physical splice is a satisfying earlier step. It distinguishes feasibility from optimization and separates an example route (an upper bound) from an obstruction (a lower bound). The upper activity turns parity into an algorithm for a small route-inspection problem.

Cornell CS 2800, Spring 2017, \emph{Lecture 36: Paths}, gives the even-degree theorem and a cycle-splicing proof. \href{https://www.cs.cornell.edu/courses/cs2800/2017sp/lectures/lec36-graphs2.html}{Official course notes}. Our maps, graded questions, and counter-based launch are original adaptations.
\Subhead{A real research connection}
Edmonds and Johnson, \emph{Matching, Euler tours and the Chinese postman}, \emph{Mathematical Programming} 5 (1973), pp. 88--124, studies the weighted route-inspection problem using matching theory. The introduction describes the problem and its structure; \S4 concerns a matching algorithm and \S5 Euler tours. \href{https://research.ibm.com/publications/matching-euler-tours-and-the-chinese-postman}{IBM publication record}; \href{https://web.eecs.umich.edu/~pettie/matching/Edmonds-Johnson-chinese-postman.pdf}{Paper hosted by University of Michigan}.

The extra page performs an exact four-odd-vertex instance of the same reduction. It does not teach the general blossom algorithm, polyhedral proof, or directed/mixed variants. For $2k$ odd vertices, exhaustive pairings number $(2k-1)(2k-3)\cdots1$; efficient weighted matching makes large instances practical. The research is a historical solved result, not a claimed present-day open problem.
\Subhead{Local sources and returners}
Rozhkovskaya, \emph{Math Circles for Elementary School Students}, M.4.1 ``Make your own problem,'' EPUB \texttt{part0031.xhtml}, inspires making under constraints. Its animal/figure/number prompt is different; our connected-map/parity specifications are our design. Lesson 3 ``At the lesson'' supports child attempts and explanations before an organizing display.

\emph{Math Circle by the Bay}, preface printed pp. viii--x, supports adjustable depth and extra challenges. Our timetable and two-adult allocation are proposals for this roster.

Lowell Handout 10's password-window puzzle is related to directed Euler/de Bruijn constructions; do not call graph traversal wholly new. Here the representation is an undirected bridge map, and the claims concern odd-degree obstructions and minimum repeated travel. Week 1's flip graph concerned states and distance; here edges are roads that must all be serviced. Generic old workshop starters remain optional future reserves.
\Subhead{Use history and checks}
IDs F10-K/M/U/X-v2. Record maps, repairs, exact routes, whether optimality was justified, and child-designed problems with their solutions. \texttt{verify.py} independently searches states (current vertex, set of roads already visited) and confirms closed/open/weighted optima 8, 7, 23. It also checks printed routes. The mathematical certificates above explain why the results hold.

\end{document}
